\documentclass[11pt]{article}
\usepackage[utf8]{inputenc}
\usepackage[T1]{fontenc}
\usepackage{amsmath, amssymb, amsthm}
\usepackage{geometry}
\geometry{a4paper, margin=2.5cm}
\usepackage{bm}

\title{Die Vektorableitung in der Geometrischen Algebra \\ nach David Hestenes}
\author{Zusammenfassung der wichtigsten Erkenntnisse}
\date{}

\begin{document}

\maketitle

\section{Grundlegende Setting}

Sei \(V\) ein endlichdimensionaler Vektorraum mit einer nicht-ausgearteten Metrik und \(\mathcal G(V)\) die zugehörige \textbf{multivektorielle Algebra} (Geometric Algebra). Wir betrachten eine Funktion

\[
f: V \longrightarrow \mathcal G(V),
\]

die also jedem Vektor \(x \in V\) einen Multivektor \(f(x) \in \mathcal G(V)\) zuordnet.

\section{Der Vektor-Differentialoperator}

Sei \(\{e_i\}\) eine Basis von \(V\) und \(\{e^i\}\) die dazu \textbf{reziproke Basis}, definiert durch

\[
\boxed{e^i \cdot e_j = \delta^i_j}.
\]

Der Vektor-Differentialoperator wird definiert als

\[
\boxed{\partial := e^i \, \partial_i},
\qquad\text{mit}\qquad
\partial_i := \frac{\partial}{\partial x^i},
\]

wobei \(x = x^i e_i\) die Koordinatendarstellung von \(x\) ist.

\begin{center}
\fbox{%
\begin{minipage}{0.9\textwidth}
\centering
\textbf{Wichtig:} Die Verwendung der reziproken Basis \(e^i\) ist essenziell für eine koordinatenunabhängige Formulierung. Bei orthonormaler Basis gilt \(e^i = e_i\).
\end{minipage}
}
\end{center}

\section{Richtungsableitung und Punktprodukt mit \(\partial\)}

Die \textbf{Richtungsableitung} von \(f\) in Richtung \(a \in V\) wird definiert durch

\[
\boxed{
(a \cdot \partial) f(x)
:=
\lim_{\varepsilon \to 0}
\frac{f(x + \varepsilon a) - f(x)}{\varepsilon}
}.
\]

Dies ist die natürliche Verallgemeinerung der Richtungsableitung auf multivektorwertige Funktionen.

\subsection{Herleitung der Komponentendarstellung}

Mit \(a = a^i e_i\) erhalten wir:

\[
a \cdot \partial
=
(a^i e_i) \cdot (e^j \partial_j)
=
a^i (e_i \cdot e^j) \partial_j
=
a^i \delta_i^j \partial_j
=
a^j \partial_j.
\]

Also gilt die Operatoridentität

\[
\boxed{
a \cdot \partial = a^i \partial_i
}.
\]

Somit ist

\[
\boxed{
(a \cdot \partial) f
=
a^i \partial_i f
}.
\]

\begin{center}
\fbox{%
\begin{minipage}{0.9\textwidth}
\centering
\textbf{Kernaussage:} Die Schreibweise \(a \cdot \partial\) ist als \emph{Richtungsableitungsoperator} zu verstehen. Die Komponentenformel folgt aus der Linearität der Richtungsableitung bezüglich \(a\) und der reziproken Basiseigenschaft.
\end{minipage}
}
\end{center}

\section{Die partielle Ableitung \(\partial_i\) auf multivektorwertigen Funktionen}

Ein Multivektor \(f(x) \in \mathcal G(V)\) kann bezüglich einer konstanten Basis \(\{E_A\}\) von \(\mathcal G(V)\) entwickelt werden:

\[
\boxed{
f(x) = \sum_A f^A(x) \, E_A,
\qquad f^A(x) \in \mathbb R.
}
\]

Dann ist die partielle Ableitung definiert durch

\[
\boxed{
\partial_i f(x)
=
\sum_A \partial_i f^A(x) \, E_A
=
\sum_A \frac{\partial f^A}{\partial x^i}(x) \, E_A.
}
\]

\begin{center}
\fbox{%
\begin{minipage}{0.9\textwidth}
\centering
\textbf{Wesentliche Erkenntnis:} Die partielle Ableitung \(\partial_i\) wirkt \emph{nur} auf die skalaren Koeffizienten \(f^A(x)\). Der geometrische Charakter (der Grad) des Multivektors bleibt erhalten.
\end{minipage}
}
\end{center}

\section{Algebraische Eigenschaften von \(\partial_x\) als Vektor}

Hestenes postuliert, dass \(\partial_x\) die \textbf{algebraischen Eigenschaften eines Vektors} in \(\mathcal G^1(V)\) besitzt. Daraus ergeben sich zwei fundamentale Operationen:

\subsection{Das äußere Produkt mit dem Pseudoskalar}

Sei \(I\) der \textbf{Pseudoskalar} des Raumes (das orientierte Volumenelement, maximaler Grad \(n\)). Dann gilt:

\[
\boxed{
I \wedge \partial_x = 0
}
\]

\textbf{Begründung:} \(I\) hat Grad \(n\), \(\partial_x\) hat Grad \(1\). Das äußere Produkt hätte Grad \(n+1\), was in einer \(n\)-dimensionalen Algebra verschwindet.

\subsection{Folgerung: Geometrisches Produkt wird zum inneren Produkt}

Da in der Geometrischen Algebra für ein geometrisches Produkt gilt:

\[
I \partial_x = I \cdot \partial_x + I \wedge \partial_x,
\]

und der äußere Term verschwindet, folgt:

\[
\boxed{
I \partial_x = I \cdot \partial_x
}
\]

\begin{center}
\fbox{%
\begin{minipage}{0.9\textwidth}
\centering
\textbf{Interpretation:} Der Pseudoskalar \(I\) und der Vektoroperator \(\partial_x\) kommutieren im Sinne des inneren Produkts. Dies ist die algebraische Grundlage für die spätere Identifikation von Divergenz und Rotation.
\end{minipage}
}
\end{center}

\section{Die koordinatenfreie Definition von \(\partial_x\) -- Der "Hammer"}

Der entscheidende Schritt in Hestenes' Formalismus ist die folgende Identität:

\[
\boxed{
\partial_x = P(\partial_x) = \sum_k a^k \, (a_k \cdot \partial_x)
}
\]

wobei \(\{a_k\}\) eine \textbf{beliebige Basis} von \(V\) ist, \(\{a^k\}\) die dazugehörige reziproke Basis, und

\[
P: \mathcal G(V) \longrightarrow \mathcal G^1(V)
\]

der \textbf{Projektionsoperator} auf den Vektorraum (den Grade-1-Teil) ist.

\begin{center}
\fbox{%
\begin{minipage}{0.9\textwidth}
\centering
\textbf{Das ist der "Hammer"!} Diese Gleichung definiert \(\partial_x\) \emph{ohne} Koordinaten, rein algebraisch, und zeigt, dass \(\partial_x\) ein echter Vektor im Sinne der Geometrischen Algebra ist.
\end{minipage}
}
\end{center}

\subsection{Bedeutung der Gleichung}

\begin{enumerate}
    \item \(\partial_x = P(\partial_x)\) bedeutet: \(\partial_x\) hat \emph{keine} Anteile höherer Grade (Bivektoren, Trivektoren usw.). Es ist ein reiner Vektoroperator.
    \item Die Summe \(\sum_k a^k (a_k \cdot \partial_x)\) ist die \textbf{Spektralzerlegung} des Identitätsoperators auf \(V\), angewendet auf den Operator \(\partial_x\).
    \item Die Gleichung gilt für \textbf{jede} Wahl der Basis \(\{a_k\}\) – das ist die koordinatenunabhängige Essenz.
\end{enumerate}

\subsection{Spezialfall: Feste Basis}

Wählt man \(a_k = e_k\) (eine feste Basis), so erhält man:

\[
\partial_x = \sum_k e^k (e_k \cdot \partial_x) = \sum_k e^k \partial_k,
\]

weil \(e_k \cdot \partial_x = \partial_k\) ist (wie in Abschnitt 3 gezeigt).

\section{Die grundlegende Ableitung: \(\partial_x x = n\)}

Die Anwendung der Vektorableitung auf die Identitätsfunktion \(F(x) = x\) liefert ein fundamentales Ergebnis:

\[
\boxed{
\partial_x x = n
}
\]

wobei \(n = \dim(V)\) die Dimension des Raumes ist.

\subsection{Beweis mit der Basis-Darstellung}

Mit dem "Hammer" aus Abschnitt 6:

\[
\partial_x x = \sum_k a^k (a_k \cdot \partial_x x) = \sum_k a^k a_k = n.
\]

Die letzte Gleichheit \(\sum_k a^k a_k = n\) folgt, weil der Skalaranteil \(n\) ist und der Bivektoranteil wegen der Symmetrie der Metrik verschwindet.

\subsection{Beweis über Divergenz und Rotation}

Alternativ:

\[
\partial_x x = \nabla \cdot x + \nabla \wedge x = n + 0 = n.
\]

\section{Die Produktregel und die Nicht-Kommutativität von \(\partial_x\)}

Die Operation mit \(\partial_x\) ist \textbf{nicht kommutativ}. Das ist der entscheidende Unterschied zur gewöhnlichen Vektoranalysis.

\subsection{Herleitung der Produktregel}

Für zwei multivektorwertige Funktionen \(F(x)\) und \(G(x)\) gilt:

\[
\partial (FG) = e^i \partial_i (FG) = e^i (\partial_i F)G + e^i F(\partial_i G).
\]

Der erste Term ist einfach \((\partial F)G\). Der zweite Term \emph{kann nicht} als \(F(\partial G)\) geschrieben werden, weil \(e^i\) und \(F\) nicht vertauscht werden dürfen!

\subsection{Die Überpunkt-Notation (Overdot Notation)}

Hestenes führt die \textbf{Überpunkt-Notation} ein, um die Wirkungsweise des Operators präzise zu kennzeichnen:

\[
\boxed{
\dot{\partial} F \dot{G} := e^i F (\partial_i G)
}
\]

Die Punkte bedeuten:
- Der Punkt über \(\partial\) zeigt: "Hier ist der Ableitungsoperator, der wirkt."
- Der Punkt über \(G\) zeigt: "Das ist die Funktion, auf die \(\partial\) wirkt."
- \textbf{Wichtig:} Die algebraische Reihenfolge der Faktoren bleibt erhalten!

Damit lautet die Produktregel:

\[
\boxed{
\partial (FG) = (\partial F)G + \dot{\partial} F \dot{G}
}
\]

oder vollständig mit Punkten:

\[
\boxed{
\partial (FG) = \dot{\partial} \dot{F} G + \dot{\partial} F \dot{G}
}
\]

\subsection{Beispiel für die Nicht-Kommutativität}

Sei \(B = e_1 e_2\) ein konstanter Bivektor im \(\mathbb R^2\). Dann gilt:

\[
\partial (B x) = \sum_{i=1}^2 e_i B e_i = -2B,
\]

während

\[
B \partial x = B \cdot 2 = 2B.
\]

Also:

\[
\boxed{
\partial (B x) = -2B \quad \neq \quad B \partial x = 2B
}
\]

\begin{center}
\fbox{%
\begin{minipage}{0.9\textwidth}
\centering
\textbf{Kernaussage:} Die Nicht-Kommutativität von \(\partial_x\) ist kein Fehler, sondern die eigentliche Stärke des Formalismus. Sie kodiert die geometrische Struktur des Raumes in den Ableitungsregeln.
\end{minipage}
}
\end{center}

\section{Zusammenhang mit der Fréchet-Ableitung}

Die Richtungsableitung von \(f\) in Richtung \(a\) ist das Fréchet-Differential \(Df(x)\) angewendet auf \(a\):

\[
\boxed{
(a \cdot \partial) f(x)
=
Df(x)[a]
}.
\]

Das Differential ist eine lineare Abbildung

\[
Df(x): V \longrightarrow \mathcal G(V).
\]

Die Komponentendarstellung lautet

\[
Df(x)[a] = a^i \partial_i f(x).
\]

\subsection{Anmerkung zur Linearität}

Die Linearität von \(Df(x)\) bezüglich \(a\) ist der Grund, warum wir die Homogenität anwenden können:

\[
Df(x)[a]
=
Df(x)[a^i e_i]
=
a^i Df(x)[e_i]
=
a^i \partial_i f(x).
\]

\begin{center}
\fbox{%
\begin{minipage}{0.9\textwidth}
\centering
\textbf{Wichtige Unterscheidung:} Der Ausdruck \(a \cdot \partial f\) ist \emph{nicht} dasselbe wie das geometrische Produkt \(a \cdot (\partial f)\) im Allgemeinen! Die Notation ist als Richtungsableitung zu lesen, nicht als Punktprodukt eines Vektors mit einem Multivektor.
\end{minipage}
}
\end{center}

\section{Zusammenfassung der wichtigsten Formeln}

\begin{align}
\partial &= e^i \partial_i, & e^i \cdot e_j &= \delta^i_j, \\
(a \cdot \partial) f(x) &= \lim_{\varepsilon \to 0} \frac{f(x + \varepsilon a) - f(x)}{\varepsilon}, \\
(a \cdot \partial) f(x) &= a^i \partial_i f(x), \\
\partial_i f(x) &= \sum_A \partial_i f^A(x) \, E_A, \\
I \wedge \partial_x &= 0, \\
I \partial_x &= I \cdot \partial_x, \\
\partial_x &= P(\partial_x) = \sum_k a^k (a_k \cdot \partial_x), \\
\partial_x x &= n, \\
\partial (FG) &= (\partial F)G + \dot{\partial} F \dot{G}, \\
(a \cdot \partial) f(x) &= Df(x)[a].
\end{align}

\begin{center}
\fbox{%
\begin{minipage}{0.9\textwidth}
\centering
\textbf{Fazit:} Die Vektorableitung von Hestenes ist eine kompakte, koordinatenunabhängige Schreibweise für die gewöhnliche Richtungsableitung auf multivektorwertigen Funktionen. Die Verwendung der reziproken Basis, die Projektion auf den Vektorraum und die Überpunkt-Notation sind die Schlüssel zum Verständnis dieses eleganten Formalismus.
\end{minipage}
}
\end{center}

\section*{Hinweis}

Die obigen Ausführungen gelten für eine \textbf{konstante Basis} \(\{E_A\}\) von \(\mathcal G(V)\). Bei ortsabhängigen Basen muss zusätzlich die Ableitung der Basis berücksichtigt werden.

\end{document}
